\documentclass[10pt, letterpaper]{article}

% Packages:
\usepackage[
ignoreheadfoot, % set margins without considering header and footer
top=1.3 cm, % seperation between body and page edge from the top
bottom=1.3 cm, % seperation between body and page edge from the bottom
left=2.0 cm, % seperation between body and page edge from the left
right=2.0 cm, % seperation between body and page edge from the right
footskip=0.7 cm, % seperation between body and footer
% showframe % for debugging
]{geometry} % for adjusting page geometry
\usepackage{titlesec} % for customizing section titles
\usepackage{tabularx} % for making tables with fixed width columns
\usepackage{array} % tabularx requires this
\usepackage[dvipsnames]{xcolor} % for coloring text
\definecolor{primaryColor}{RGB}{0, 102, 204} % define primary color as blue
\definecolor{darkRed}{RGB}{128, 0, 32} % dark red for highlighted distinctions
\usepackage{enumitem} % for customizing lists
\usepackage{fontawesome5} % for using icons
\usepackage{amsmath} % for math
\usepackage[
pdftitle={Xinglang Zhang - Research CV},
pdfauthor={Xinglang Zhang},
pdfcreator={LaTeX with RenderCV},
colorlinks=true,
urlcolor=primaryColor
]{hyperref} % for links, metadata and bookmarks
\usepackage[pscoord]{eso-pic} % for floating text on the page
\usepackage{calc} % for calculating lengths
\usepackage{bookmark} % for bookmarks
\usepackage{lastpage} % for getting the total number of pages
\usepackage{changepage} % for one column entries (adjustwidth environment)
\usepackage{paracol} % for two and three column entries
\usepackage{ifthen} % for conditional statements
\usepackage{needspace} % for avoiding page brake right after the section title
\usepackage{iftex} % check if engine is pdflatex, xetex or luatex

% Ensure that generate pdf is machine readable/ATS parsable:
\ifPDFTeX
\input{glyphtounicode}
\pdfgentounicode=1
\usepackage[T1]{fontenc}
\usepackage[utf8]{inputenc}
\usepackage{lmodern}
\fi

\usepackage{charter}

% Compact Times New Roman italic annotations for parenthetical distinctions in Honors.
% pdfLaTeX uses its Times-compatible fallback; XeLaTeX embeds the installed font.
\ifPDFTeX
    \newcommand{\honornote}[1]{{\small\fontfamily{ptm}\selectfont\itshape (#1)}}
\else
    \usepackage{fontspec}
    % Keep the original pdfLaTeX Charter family for all resume body text.
    \AtBeginDocument{\fontencoding{T1}\fontfamily{bch}\selectfont}
    \newfontfamily\honornotefont{Times New Roman}
    \newcommand{\honornote}[1]{{\small\honornotefont\itshape (#1)}}
\fi
\newcommand{\distinction}[1]{\textcolor{darkRed}{#1}}
\newcommand{\pubnum}[1]{{\footnotesize\textbf{[#1]}}}

% Some settings:
\raggedright
\AtBeginEnvironment{adjustwidth}{\partopsep0pt} % remove space before adjustwidth environment
\pagestyle{empty} % no header or footer
\setcounter{secnumdepth}{0} % no section numbering
\setlength{\parindent}{0pt} % no indentation
\setlength{\topskip}{0pt} % no top skip
\setlength{\columnsep}{0.15cm} % set column seperation
\pagenumbering{arabic}
\pagestyle{plain}
\linespread{0.97} % compact line spacing

\titleformat{\section}{\needspace{3\baselineskip}\bfseries\large}{}{0pt}{}[\vspace{1pt}\titlerule]

\titlespacing{\section}{
    -1pt
}{
    0.22 cm
}{
    0.12 cm
}


\renewcommand\labelitemi{$\vcenter{\hbox{\small$\bullet$}}$} % custom bullet points
\newenvironment{highlights}{
    \begin{itemize}[
    topsep=0.05 cm,
    parsep=0.02 cm,
    partopsep=0pt,
    itemsep=0pt,
    leftmargin=0 cm + 10pt
    ]
}{
\end{itemize}
} % new environment for highlights

\newenvironment{highlightsforbulletentries}{
    \begin{itemize}[
    topsep=0.05 cm,
    parsep=0.02 cm,
    partopsep=0pt,
    itemsep=0pt,
    leftmargin=10pt
    ]
}{
\end{itemize}
} % new environment for highlights for bullet entries

\newenvironment{onecolentry}{
    \begin{adjustwidth}{
        0 cm + 0.00001 cm
    }{
        0 cm + 0.00001 cm
    }
}{
\end{adjustwidth}
} % new environment for one column entries

\newenvironment{twocolentry}[2][]{
    \onecolentry
    \def\secondColumn{#2}
    \setcolumnwidth{\fill, 4.5 cm}
    \begin{paracol}{2}
}{
\switchcolumn \raggedleft \secondColumn
\end{paracol}
\endonecolentry
} % new environment for two column entries

\newenvironment{threecolentry}[3][]{
    \onecolentry
    \def\thirdColumn{#3}
    \setcolumnwidth{, \fill, 4.5 cm}
    \begin{paracol}{3}
    {\raggedright #2} \switchcolumn
    }{
    \switchcolumn \raggedleft \thirdColumn
    \end{paracol}
    \endonecolentry
} % new environment for three column entries

\newenvironment{header}{
    \setlength{\topsep}{0pt}\par\kern\topsep\centering\linespread{1.5}
}{
    \par\kern\topsep
} % new environment for the header

\newcommand{\placelastupdatedtext}{% \placetextbox{<horizontal pos>}{<vertical pos>}{<stuff>}
\AddToShipoutPictureFG*{% Add <stuff> to current page foreground
\put(
\LenToUnit{\paperwidth-2 cm-0 cm+0.05cm},
\LenToUnit{\paperheight-1.0 cm}
){\vtop{{\null}\makebox[0pt][c]{
    \small\color{gray}\textit{Last updated in October 2026}\hspace{\widthof{Last updated in October 2026}}
}}}%
}%
}%

% save the original href command in a new command:
\let\hrefWithoutArrow\href

\begin{document}
\newcommand{\AND}{\unskip
\cleaders\copy\ANDbox\hskip\wd\ANDbox
\ignorespaces
}
\newsavebox\ANDbox
\sbox\ANDbox{$|$}

\begin{header}
\fontsize{25 pt}{25 pt}\selectfont Xinglang Zhang

\vspace{5 pt}

\normalsize
\mbox{\hrefWithoutArrow{mailto:normanspark@hust.edu.cn}{normanspark@hust.edu.cn}}%
\kern 5.0 pt%
\AND%
\kern 5.0 pt%
\mbox{\hrefWithoutArrow{https://github.com/DelicateNorman}{\color{blue}GitHub}}%
\kern 5.0 pt%
\AND%
\kern 5.0 pt%
\mbox{\hrefWithoutArrow{https://delicatenorman.github.io}{\color{blue}Personal Page}}%
\kern 5.0 pt%
\end{header}

\vspace{5 pt - 0.3 cm}

\section{Education}

\begin{twocolentry}{
    Sept 2024 -- Jun 2028
}
\textbf{Huazhong University of Science and Technology (HUST)}
\end{twocolentry}

\vspace{0.10 cm}

\begin{onecolentry}
    B.S. in Cyber Science and Engineering
\end{onecolentry}

\vspace{0.10 cm}

\begin{onecolentry}
    \begin{highlights}
        \item \textbf{GPA:} 4.80/5.0 \quad \textbf{Weighted Avg:} 94.51 \quad \textbf{Rank:} \textbf{\distinction{1/108}}
        \item \textbf{Selected Coursework:} Calculus A (I/II) (A+), Linear Algebra (A+), Discrete Math (A+), Data Structures (A+), Assembly Language (A+), Probability Theory (A+), Digital Circuit (A+), Database Systems (A+).
    \end{highlights}
\end{onecolentry}

\section{Research Agenda}
\begin{onecolentry}
\textbf{Reliable LLM Reasoning.} When and why does reasoning break down, and how can structured or neuro-symbolic mechanisms make it more reliable?\par
\smallskip
\textbf{Test-Time Learning \& Self-Evolving Agents.} How can deployed models learn from their own experience without gold supervision, and decide what, when, and how to learn as their state evolves?\par
\smallskip
\textbf{Human--AI Preference \& Social Intelligence.} Where does model behavior diverge from human preferences and social signals, and how can these gaps inform model design?
\end{onecolentry}

\section{Research Experience}

\newcommand{\researchtheme}[1]{\par\addvspace{0.16cm}\needspace{6\baselineskip}{\color{primaryColor}\bfseries #1}\par\nobreak\vspace{0.08cm}}
\newcommand{\researchproject}[2]{\needspace{6\baselineskip}\textbf{#1}\hfill {\small #2}\par\nobreak\vspace{0.05cm}}

% Narrative evidence: the PDF linked by publication [3] in the original CV.
% https://drive.google.com/file/d/1_pGifgLGeNbYhm97wl-YnW84gn0-pHmm/view?usp=sharing
% Preserve the linked main-result range: 13.10--24.42 pp.
\researchtheme{Adaptive Learning \& Self-Evolving Agents}
\researchproject{TRACE: Transfer-Aware Test-Time Training for Tool-Using LLMs}{Jul 2026 -- Present}
\begin{onecolentry}
\textbf{Research Intern}, BLENDER Lab, University of Illinois Urbana-Champaign (UIUC) \\
\textbf{Mentor:} \href{https://qiancheng0.github.io/}{Cheng Qian}
% \quad\textbf{Advisor:} \href{https://blender.cs.illinois.edu/hengji.html}{Heng Ji}
\begin{highlights}
\item \textbf{Question.} Test-time training can learn from self-generated tool-use practices, but a limited update budget makes experience selection essential. Formulated the question: \emph{which experience is worth learning now, given what the model has already learned?}
\item \textbf{Finding.} Characterized \emph{Transfer Value} as the benefit of learning from an experience for a target distribution. Established that this utility is \textbf{relational and state-dependent}: the same experience can help one query and harm another, and its value can reverse as the solver updates. Demonstrated cross-query transfer and non-monotonic changes in utility. In a controlled probe, \textbf{30 of 32} initially beneficial experiences became harmful after 32 updates; some became useful again later, motivating repeated selection rather than a fixed ranking.
\item \textbf{Method.} Developed \emph{TRACE} to combine \emph{Target Relevance}, \emph{Current-State Transferability}, and \emph{Marginal Complementarity}. Temporary updates and gold-free probes allocate the learning budget; complementary practices are selected, learned through persistent LoRA-SFT, and re-evaluated as the solver evolves.
\item \textbf{Evidence.} Under matched permanent update budgets, probing the updated solver reached \textbf{38.1 Exact} versus \textbf{30.4} with stale-state probes. Reordering the same selected experiences degraded performance, showing that both \emph{what} and \emph{when} to learn matter. Across four tool-use benchmarks and four main models, improved Overall Exact accuracy by \textbf{13.10--24.42 percentage points} over frozen models and consistently outperformed static selection. These findings frame effective TTT as \textbf{dynamic curriculum formation}. First-author ACL 2027 submission, under review \pubnum{3}.
\end{highlights}
\end{onecolentry}

\researchtheme{Reliable LLM Reasoning}
% Narrative evidence: https://arxiv.org/abs/2601.02902 (v2).
\researchproject{Logical Phase Transitions: Diagnosing and Mitigating Reasoning Collapse}{Aug 2025 -- Jan 2026}
\begin{onecolentry}
\textbf{Project Leader}, Intelligent Media Computing and Network Security Lab, HUST \\
\textbf{Advisor:} \href{https://scholar.google.com/citations?user=1qnuOZsAAAAJ&hl=en}{Zikai Song}
\begin{highlights}
\item Investigated where reasoning becomes unreliable as \emph{logical complexity} increases, asking whether collapse follows identifiable regimes that aggregate benchmark scores obscure.
\item \textbf{LoCM and NSA-LR.} Introduced the \emph{Logical Complexity Metric (LoCM)} and \emph{NSA-LR}, pairing natural-language statements with explicit first-order-logic representations and reasoning chains. Enabled complexity analysis beyond coarse hop-count estimates.
\item Identified \emph{Logical Phase Transitions}: accuracy remains relatively stable within complexity regimes, then drops sharply across critical intervals. Observed this pattern across open- and closed-source LLMs.
\item Developed \emph{Neuro-Symbolic Curriculum Tuning}, combining natural-language/symbolic alignment with complexity-aware training around the observed transition regions to progressively strengthen reasoning at higher complexity.
\item Across five benchmarks, obtained average accuracy gains of \textbf{+1.26} with naive prompting and \textbf{+3.95} with CoT, alongside improved generalization to unseen logical compositions. Connected failure diagnosis to curriculum design. Co-first-author paper, \textbf{\distinction{ACL 2026 Main}} \pubnum{1}.
\end{highlights}
\end{onecolentry}

% Narrative evidence: https://arxiv.org/abs/2509.24765 and https://arxiv.org/abs/2607.21933.
\newpage
\section{Research Experience (continued)}
\researchtheme{Reliable LLM Reasoning (continued)}
\researchproject{From LogicAgent to HexLogicAgent: Structuring Meaning Before Deduction}{}
\begin{onecolentry}
\textbf{Research Assistant}, Intelligent Media Computing and Network Security Lab, HUST \\
\textbf{Advisor:} \href{https://scholar.google.com/citations?user=1qnuOZsAAAAJ&hl=en}{Zikai Song} \\
{\small LogicAgent: Feb -- Aug 2025\quad HexLogicAgent: Mar -- Jul 2026}
\begin{highlights}
\item Investigated how \emph{semantic complexity} interacts with \emph{logical depth}: valid deduction depends on first interpreting abstract, ambiguous, or conflicting propositions correctly.
\item \textbf{LogicAgent.} Contributed to a semiotic-square-guided framework that structures alternative meanings before automated deduction and reflective verification. Introduced \emph{RepublicQA} to jointly evaluate semantic and logical difficulty; achieved \textbf{+6.25\%} average improvement on RepublicQA and \textbf{+7.05\%} across four mainstream benchmarks. Co-authored \textbf{\distinction{ACL 2026 Main / Oral}} paper \pubnum{2}.
\item The square leaves mixed and partially compatible meanings underrepresented. Identified incomplete semantic representations as a major source of failures despite apparently valid deduction.
\item \textbf{HexLogicAgent.} Extended the semantic structure from a square to a logical hexagon, explicitly representing uniform and mixed cases before deduction and verification. Developed a formal account of semantic opposition and completeness under stated first-order-logic assumptions.
\item HexLogicAgent improved reliability across five benchmarks and delayed degradation as logical complexity increased. Established \textbf{semantic organization before inference} as a focus of this research line. Co-first-author \emph{Artificial Intelligence} submission \pubnum{5}.
\end{highlights}
\end{onecolentry}

\researchtheme{Human--AI Preference \& Social Intelligence}
% Narrative evidence: https://arxiv.org/abs/2609.18282.
% Review scores and submission status are preserved from the original CV.
\researchproject{OMRA: Diagnosing and Reducing Synthetic--Real Preference Misalignment}{Jan -- May 2026}
\begin{onecolentry}
\textbf{Project Leader}, Intelligent Media Computing and Network Security Lab, HUST \\
\textbf{Advisor:} \href{https://scholar.google.com/citations?user=1qnuOZsAAAAJ&hl=en}{Zikai Song}
\begin{highlights}
\item Asked whether the qualities LLMs associate with high engagement match real user responses. Studied \textbf{1.17 million answers} to \textbf{25,978 questions} on Zhihu, Quora, and Reddit; introduced \emph{Ontological Preference Measurement} along logic, affect, and expression dimensions.
\item Identified \emph{logic overbinding}: as target engagement increases, LLMs add explicit reasoning structure, while observed user engagement is more strongly associated with affective and expressive salience.
\item Developed \emph{Ontology-Masked Reasoning Autoencoding (OMRA)} to mask and reconstruct over-explained spans while preserving stance, factual content, and coherence.
\item Reduced marginal, geometric, and transfer gaps by \textbf{54.4\% on average} across four LLM families. The results motivate grounding preference alignment in observed social signals. First-author NAACL 2027 submission \pubnum{4}; all three reviewers gave \textbf{OA 3.5/5}, with a reported top-8\% ranking.
\end{highlights}
\end{onecolentry}

\begin{onecolentry}
\textbf{Related collaborative work.}
\begin{highlights}
\item Co-authored \emph{IntervenSim} \pubnum{7}, modeling the co-evolution of interventions and crowd opinion; co-authored the \emph{Social Intelligence Modeling} survey \pubnum{6}, connecting social perception with social simulation.
\end{highlights}
\end{onecolentry}

\section{Publications}
{\small * Equal contribution.}\par\smallskip

\begin{onecolentry}
\pubnum{1}~\textbf{Logical Phase Transitions: Understanding Collapse in LLM Logical Reasoning} \\
\textit{\textbf{Xinglang Zhang*}, Yunyao Zhang*, et al.} | \textbf{\distinction{ACL 2026} (Main; co-first author)} |
\href{https://arxiv.org/abs/2601.02902}{\color{blue}[ArXiv]} |
\href{https://github.com/AI4SS/Logical-Phase-Transitions}{\color{blue}[Code]}
\end{onecolentry}
\vspace{0.08 cm}

\begin{onecolentry}
\pubnum{2}~\textbf{Semantic-Aware Logical Reasoning via a Semiotic Framework} \\
\textit{Yunyao Zhang, \textbf{Xinglang Zhang}, et al.} | \textbf{\distinction{ACL 2026} (Main; \distinction{Oral})} |
\href{https://arxiv.org/abs/2509.24765}{\color{blue}[ArXiv]} |
\href{https://github.com/AI4SS/Logic-Agent}{\color{blue}[Code]}
\end{onecolentry}
\vspace{0.08 cm}

\begin{onecolentry}
\pubnum{3}~\textbf{TRACE: Transfer-Aware Test-Time Training for Tool-Using LLMs} \\
\textit{\textbf{Xinglang Zhang}, Cheng Qian} |
\href{https://drive.google.com/file/d/1_pGifgLGeNbYhm97wl-YnW84gn0-pHmm/view?usp=sharing}{\color{blue}[Paper]} |
\textbf{ACL 2027, under review}
\end{onecolentry}
\vspace{0.08 cm}

\begin{onecolentry}
\pubnum{4}~{\small\textbf{Too Good to Be Real? Diagnosing and Reducing the Gap Between AI Preference and Real User Engagement}} \\
\textit{\textbf{Xinglang Zhang}, Yuanmeng Xiang, et al.} |
\href{https://arxiv.org/abs/2609.18282}{\color{blue}[ArXiv]} |
\textbf{NAACL 2027 submission}
\end{onecolentry}
\vspace{0.08 cm}

\begin{onecolentry}
\pubnum{5}~\textbf{Semiotic Logical Hexagon Theory for LLM Logical Reasoning} \\
\textit{Yunyao Zhang*, \textbf{Xinglang Zhang*}, et al.} |
\href{https://arxiv.org/abs/2607.21933}{\color{blue}[ArXiv]} |
\href{https://github.com/AI4SS/Logic-Agent}{\color{blue}[Code]} |
\textbf{Co-first author; Artificial Intelligence submission}
\end{onecolentry}
\vspace{0.08 cm}

\begin{onecolentry}
\pubnum{6}~\textbf{Social Intelligence Modeling: A Comprehensive Survey from Social Perception to Social Simulation} \\
\textit{Zikai Song, Xiajie Li, Yunyao Zhang, \textbf{Xinglang Zhang}, et al.} |
\href{https://www.researchgate.net/publication/411123257_Social_Intelligence_Modeling_A_Comprehensive_Survey_from_Social_Perception_to_Social_Simulation}{\color{blue}[Paper]} |
\href{https://github.com/sait-crypto/Awesome-Social-Intelligence-Modeling-System}{\color{blue}[GitHub]}
\end{onecolentry}
\vspace{0.08 cm}

\begin{onecolentry}
\pubnum{7}~\textbf{IntervenSim: Intervention-Aware Social Network Simulation for Opinion Dynamics} \\
\textit{Yunyao Zhang, Zuocheng Ying, \textbf{Xinglang Zhang}} |
\href{https://arxiv.org/abs/2604.06600}{\color{blue}[ArXiv]} |
\href{https://github.com/AI4SS/IntervenSim}{\color{blue}[Code]}
\end{onecolentry}
\vspace{0.08 cm}

\begin{onecolentry}
\pubnum{8}~\textbf{EarthVerse: Benchmarking Scientific Agents Across Dynamic Earth Systems and Natural Hazards} \\
\textit{Zhiqing Cui, Xinxiang Yin, Yihong Tang, \textbf{Xinglang Zhang}, et al.} |
\href{https://cuizhiq.github.io/EarthVerse/}{\color{blue}[Project Page]} |
\href{https://arxiv.org/abs/2608.23525}{\color{blue}[Paper]}
\end{onecolentry}

\newpage
\section{Selected Funded Projects}

\begin{onecolentry}
\begin{itemize}[leftmargin=10pt, topsep=0.10cm, itemsep=0.22cm, parsep=0pt]
\item \textbf{HUST Undergraduate Natural Science Fund, 2026} \\
\textit{Trustworthy Enhancement for Complex Reasoning} \\
Project Lead \quad | \quad \textbf{CNY 50,000 (approx. US\$7,450)}
\item \textbf{College Student Innovation and Entrepreneurship Program (University Level)} \\
\textit{LLM Logical Reasoning through Philosophical Paradigms} \\
Project Lead \quad | \quad \textbf{CNY 4,000 (approx. US\$600)}
\item \textbf{College Student Innovation and Entrepreneurship Program (Hubei Provincial Level)} \\
\textit{Trustworthy Enhancement of LLMs via a Neuro-Symbolic Framework} \\
Project Lead \quad | \quad \textbf{CNY 20,000 (approx. US\$2,980)}
\end{itemize}
\end{onecolentry}

\vspace{0.12cm}
\section{Honors}

% Keep distinctions inline, with extra space between awards.
\newcommand{\honorentry}[3]{\item \textbf{#1} \honornote{#2}\hfill #3}
\begin{onecolentry}
\begin{itemize}[leftmargin=10pt, topsep=0.10cm, itemsep=0.12cm, parsep=0pt]
\honorentry{Outstanding Student, Qiming College}{Highest Honor for HUST Undergraduates; \distinction{Top 1.85\%}}{2025 -- 2027}
\honorentry{Jin Yugang Lan-Ying Scholarship}{\distinction{Top 1\%}}{2025 -- 2026}
\honorentry{National Scholarship}{Highest Honor for Chinese Undergraduates; \distinction{Top 0.3\% Nationwide}}{2025 -- 2026}
\honorentry{Triple-A ("Sanhao") Student Scholarship}{\distinction{Top 5.5\%}}{2025 -- 2026}
\honorentry{National Scholarship}{Highest Honor for Chinese Undergraduates; \distinction{Top 0.3\% Nationwide}}{2024 -- 2025}
\honorentry{Triple-A ("Sanhao") Student Scholarship}{\distinction{Top 5.5\%}}{2024 -- 2025}
\honorentry{Outstanding Youth League Member}{\distinction{Top 4.1\%}}{2024 -- 2025}
\honorentry{Outstanding Youth League Cadre}{\distinction{Top 4.1\%}}{2024 -- 2025}
\honorentry{Academic Excellence Scholarship \& Self-improvement Scholarship}{\distinction{Top 10\%}}{2024 -- 2025}
\end{itemize}
\end{onecolentry}

\vspace{0.12cm}
\section{Service}

\begin{onecolentry}
\textbf{Reviewer, ACL ARR} \hfill January, May, and August 2026
\end{onecolentry}

\section{Academic Activities}

\begin{onecolentry}
\textbf{ACL 2026 Poster Presenter and Oral Presenter} \hfill July 1 -- 8, 2026
\end{onecolentry}

\begin{onecolentry}
\textbf{Selected Rising Star Poster Presenter}, X-AGI Conference \hfill Oct. 17 -- 18, 2026 \textit{(upcoming)}
\end{onecolentry}

\section{Skills}

\begin{onecolentry}
\textbf{ML \& Research Stack:} Python, C/C++, PyTorch, HuggingFace, LoRA, SFT, RLHF, vLLM
\end{onecolentry}

\vspace{0.15 cm}

\begin{onecolentry}
\textbf{Mathematical Foundations:} Linear Algebra, Probability Theory, Calculus, Discrete Mathematics
\end{onecolentry}

\vspace{0.15 cm}

\begin{onecolentry}
\textbf{Languages:} English (TOEFL 5.5/6; W6 R6 L5.5 S5; CET-4: 609; CET-6: 611), Chinese (Native)
\end{onecolentry}

\section{Leadership Experience}

\begin{onecolentry}
    \begin{highlights}
        \item \textbf{Class Representative, Qiming Experimental Class 2401} \hfill Sept 2024 -- Present
        \begin{highlightsforbulletentries}
            \item Organized academic events, managed administrative tasks, and mentored students to enhance class cohesion.
        \end{highlightsforbulletentries}

        \item \textbf{Head of Organization Department, School of Cyber Science and Engineering} \hfill Sept 2025 -- Present
        \begin{highlightsforbulletentries}
            \item Managed administrative and organizational affairs for the Communist Youth League.
        \end{highlightsforbulletentries}

        \item \textbf{Head of Lesson Planning, Peer Lecturer Group} \hfill Apr 2025 -- Present
        \begin{highlightsforbulletentries}
            \item Fostered academic excellence across the college through peer mentoring and curriculum planning.
        \end{highlightsforbulletentries}
    \end{highlights}
\end{onecolentry}

\end{document}
